% !TeX root = ../papers/nc-p3.tex
\section{正規一次直線列の四窓からの有限認識と整三次商}
\label{sec:normal-line-recognition}
前節の分類枝に属するかどうかを、直線列の指定も正則性の事前認証も使わずに調べる。
まず三生成商に必要な古典的な幾何的正則性判定を、本稿の関係空間の規約で示す。

\begin{lemma}\label{lem:ternary-finite-graph-test}
$S=T(V)/(R)$, $\dim V=\dim R=3$ とする。標準三次元Koszul左右AS正則性の必要十分条件は次である。
\begin{enumerate}
\item $(R\otimes V)\cap(V\otimes R)$ に、左右の一文字縮約がともに $R$ を張る元 $w$ がある。
\item $R$ の三つの双一次式の左右評価行列は射影平面の各点で階数2以上であり、
両行列式はともに恒等零か、同じ三次曲線をスキームとして定める。
\end{enumerate}
\end{lemma}
\begin{proof}
必要性は三次双対のFrobenius対合と補題\ref{lem:ternary-point-graph-growth}である。
十分性を示す。二次小行列が単元の開集合で核直線を取ると、三式の共通零点スキームは
両側の行列式軌跡上のグラフであり、両射影は同型である。平面全体ならその自己同型は
射影線型で、関係空間は平面グラフの双一次式全体となる。これは可換多項式環の次数twistである。

三次曲線 $E$（非被約でもよい）の場合、グラフ移動を $\sigma$、$L_j=(\sigma^j)^*\mathcal O_E(1)$、
$L_{[i,j)}=\bigotimes_{i\le h<j}L_h$ と置く。グラフは $\mathbb P^2\times\mathbb P^2$ で
余次元3の三式による完全交叉である。Koszul層解消を $(1,1)$ で捩ると、
$\mathcal O(-1,-1),\mathcal O(-2,-2)$ の全コホモロジーが消えることから
二文字の評価は $H^0(L_{[0,2)})$ へ全射で核はちょうど $R$ である。
非空な積 $L_{[i,j)}$ では $H^1=0$ である。実際 $L_i$ では平面三次式の完全列による。
新しい因子の基点自由な切断から全随伴点を避けるものを選ぶと、それによる乗法は単射で
余核の台が零次元なのでこの消滅は伝わる。同じ完全列と各 $L_j$ の次数3により
$h^0(L_{[0,n)})=3n$ ($n>0$) となる。

$D=\bigoplus_{n\ge0}H^0(L_{[0,n)})$ に $\sigma$ を用いた切断の積を入れる。
定理\ref{thm:maximal-tangent-complete}に用いた滑らかな場合と同じく、生成元・関係・
$w$ の複体を $E$ 上で評価すると $n\ge3$ で
\[
 0\to L_{[3,n)}\to R\otimes L_{[2,n)}
 \to V\otimes L_{[1,n)}\to L_{[0,n)}\to0
\]
が完全となる。合成は関係と三次交叉で零で、幾何学的各点の階数は $1,2,1$ である。
階数の小行列によって複体は局所的に分裂するので非被約点でも成立する。
大域切断では $n=3$ における $H^1(\mathcal O_E)=k$ だけが
$V\otimes H^0(L_{[1,3)})$ の位置のホモロジーを与える。他の正の次数ではここは完全で、
$R\otimes H^0(L_{[2,n)})$ の核は常に $w$ の像である。よって $D$ は一次生成である。
右単純加群の極小解消における最初の二微分は $V,R$ からなる。第三の微分の像には
$w$ が含まれるので、$D$ をテンソルした今の計算により
$\operatorname{Tor}_1^S(k,D)=k(-3)$, $\operatorname{Tor}_2^S(k,D)=0$ である。
従って左加群 $D$ の極小解消は $0\to S(-3)\to S\to D\to0$ となり、
左端は三次元 $g$ の右乗法で単射である。$H_S(z)=(1-z)^{-3}$ が従う。
$D$ に左socleはない。基点自由な $L_0$ の非零因子切断を掛ける写像は、
$\sigma$ による引戻しを合わせても単射である。従って $S$ の左socle元は $a g$ と書け、
$g$ による単射性から $a$ も左socleに属する。次数降下で零を得る。反対側でも同じである。
$w$ の端の縮約が全階数なので通常の長さ3のKoszul複体の左端は単射、
生成元と関係で右側は完全である。Hilbert級数のEuler等式から残る中間項も完全、
双対の相補積は両端の縮約から完全である。
\end{proof}

\begin{theorem}\label{thm:finite-normal-line-recognition}
標準四テンソル表示の二隣接切断が階数6とする。以下の有限条件は、
この窓が標準Koszul左右AS正則で正則正規一次直線列を持つことと同値である。
関係を $R_i\subset V_i\otimes V_{i+1}$ と書き、$i$ は法4とする。
\begin{enumerate}
\item 各 $V_i$ に非零 $t_i$ を取り、
$(V_i\otimes V_{i+1})/R_i$ 内の $t_i V_{i+1}$ と $V_i t_{i+1}$ が
同じ四次元空間になることを課す。
\item $x t_{i+1}=t_i g_i(x)$ で一意に定まる可逆行列について
$(g_i\otimes g_{i+1})R_i=R_{i+1}$ を四位置で検査する。
\item $g_0$ で最初の関係を一つの四次元空間の二乗に移し、$t=t_0$ で割る。
三生成商の関係次元3、補題\ref{lem:ternary-finite-graph-test}の二条件、
および持上げの $\alpha,\beta$ に対する定理\ref{thm:regular-normal-linear}の三条件を検査する。
\end{enumerate}
全ての正規直線列はこの四直線の解として現れる。従って二つの認証済み対象の
直線を指定しない同型は、一方の印を一つ固定し、他方の四直線の解と
定理\ref{thm:regular-normal-linear}の二式を同時に解いて判定できる。
\end{theorem}
\begin{proof}
四直線が得られれば $g_i$ は一移動の同型を与え、$g_i(t_i)=t_{i+1}$ である。
これに沿って全位置の基底を運べば、最初の六関係の反復で $t$ が中心になる。
三つの独立な中心関係以外の持上げを取って前節の定理を使うと十分性が従う。
判定は有限な多項式の等式・非消滅だけでよい。例えば射影点での階数1以下の不在は
各アフィンチャート上の小行列式のイデアルで検査できる。

必要性について、任意の印は前節で得る中心正則元 $t$ の形にすることができるが、
四個で十分とするにはFrobenius標準四移動がこの印を保つことを示す必要がある。
$A\to S=A/(t)$ の自然な二次双対写像 $S^!\to A^!$ は単射である。
実際、前節の変換環Tor系列では両段が線型で全微分が零であるから下段への辺準同型
$\operatorname{Tor}^A_i(k,k)\to\operatorname{Tor}^S_i(k,k)$ は全射である。
双対を取りKoszul双対との自然な同定を使えばよい。特に $S^!_3$ の生成元には
非零像 $u\in A^!_3$ がある。$t^\perp=S^!_1$ は左右とも $u$ を消し、
$A^!$ のFrobenius完全対合から両零化空間は正確に $t^\perp$ である。
$A^!$ のNakayama行列はこの超平面を保つ。実際、
$\operatorname{tr}(v u)=\operatorname{tr}(u\nu(v))$ を使えばよい。
よって逆双対で得る標準四移動は $kt$ を保つ。
代表元の倍率は自由なので厳密四周期の線として記録できる。
これで多くの印を持つ場合にも全印が有限な四直線検査から漏れない。
\end{proof}

\begin{proposition}\label{prop:integral-singular-central}
整な節点・尖点三次曲線のグラフを点対応に持つ商でも、前節の一次直線列は一意である。
持上げ変更 $e\mapsto D(e)$ は単射で、補正の重み1、2、3の方程式と有限射影同型検査がそのまま使える。
節点商の射影自己同型候補は高々六つ、尖点商では恒等のみである。
\end{proposition}
\begin{proof}
滑らかな場合の零因子の証明には超平面切断が有効Cartier因子であることだけで足り、
整三次曲線ではこれが成立する。$\sigma^*L\not\simeq L$ も同じく二次式の制限から従う。
$D(e)=0$ なら対角との交叉の接次元が2となるので $e$ は特異点で、
しかもグラフの接空間が対角内にある。従って $d\sigma_e=1$ である。

節点の場合、正規化で二枝を $0,\infty$ に取ると自己同型は $z\mapsto a z$ または $a/z$。
後者は枝を交換し、前者の二接線倍率は $a,a^{-1}$ であるから恒等微分を持つのは恒等だけである。
射影変換に延びるものは、節点の表示 $[z:z^2:1+z^3]$ の三次切断空間を保つので
$a^3=1$ の高々六つである。尖点では唯一の滑らかな変曲点を $u=0$ に取って
正規化の切断空間を $\langle1,u,u^3\rangle$ と書ける。自己同型は $u\mapsto a u+b$。
これが切断空間を保つ条件は $b=0$ である。従って $\sigma$ は $b\ne0$ を持ち、
これと可換する射影倍率 $u\mapsto d u$ は $d=1$ だけである。
$s=1/u$ として尖点の極大イデアルをその平方で割ると基底は $s^2,s^3$ である。
引戻しで $s^2\mapsto a^{-2}s^2-2ba^{-3}s^3$, $s^3\mapsto a^{-3}s^3$ となり、
恒等微分を持つには $a=1,b=0$ が必要である。両曲線で恒等移動は偏極を保つため除外され、
$D(e)=0$ から矛盾を得た。
\end{proof}
